\begin{frame}[label=inicio]{Introducción}

\centering
{\color{Cian}\bfseries Problema Industrial}\par
\vspace{0.05cm}
{\Large\bfseries Estabilidad Térmica vs. Eficiencia Catalítica}\par
\vspace{0.03cm}
{\scriptsize\color{TextoSuave}Biocatalizadores industriales: durar más sin convertir peor}

\vspace{0.14cm}
\begin{alertblock}{Pregunta Central}
\centering
\small\bfseries
¿Aumentar la termoestabilidad reduce necesariamente la eficiencia catalítica?
\end{alertblock}

\vspace{0.09cm}
\begin{columns}[T,onlytextwidth]
\column{0.31\textwidth}
\tarjeta{Estabilidad}{\centering Vida operativa\\$t_{50}$ + actividad residual}
\column{0.31\textwidth}
\tarjeta{Eficiencia}{\centering Conversión\\$V_{\max}/K_m$}
\column{0.31\textwidth}
\tarjeta{Producción}{\centering Resultado final\\producto acumulado}
\end{columns}

\vspace{0.10cm}
\begin{beamercolorbox}[wd=0.86\textwidth,sep=0.10cm,rounded=true,center]{block body}
\small\bfseries\color{CianClaro}
Estabilidad $+$ Eficiencia $+$ Operación $\longrightarrow$ Producción Útil
\end{beamercolorbox}

\vfill
\fuente{Hou y cols. (2023); Siddiqui (2017).}
\end{frame}
